\subsection{Partially isometric linear mappings}

DEFINITION 2. --- Let E and F be two hilbertian spaces and \(u \in \mathcal{L}(E; F)\) . The orthogonal of the kernel of u in E is said to be the initial subspace of u and the closure of the image of u in F is called the final subspace of u. The orthoprojector from E (resp. F) onto the initial (resp. final) subspace of u is called the initial (resp. final) orthoprojector of u.

Let \(\mathbf{P}\) be the initial subspace of \(u\) . Since \(\mathbf{E}\) is the direct sum of \(\mathbf{P}\) and of the kernel of \(u\) , we have \(u(\mathbf{P}) = u(\mathbf{E})\) .

PROPOSITION 4. --- (i) The initial (resp. final) subspace of \(u^*\) is equal to the final (resp. initial) subspace of \(u\) .

\begin{enumerate}
\def\labelenumi{(\roman{enumi})}
\setcounter{enumi}{1}
\tightlist
\item
  Suppose that \(\mathbf{E} = \mathbf{F}\) . Let \(\mathbf{M}\) be a closed vector subspace of \(\mathbf{E}\) and \(\mathbf{M}^{\circ}\) its orthogonal. The relations \(u(\mathbf{M}) \subset \mathbf{M}\) and \(u^{*}(\mathbf{M}^{\circ}) \subset \mathbf{M}^{\circ}\) are equivalent.
\end{enumerate}

Let \(\mathrm{Q} = \overline{u(\mathrm{E})}\) be the final subspace of \(u\) . The orthogonal \(\mathrm{Q}^{\circ}\) of \(\mathrm{Q}\) in \(\mathrm{F}\) consists of all vectors \(y\) such that \(\langle u(x)|y\rangle = 0\) for all \(x\in \mathrm{E}\) ; this is equivalent to: \(\langle x|u^{*}(y)\rangle = 0\) for all \(x\in \mathrm{E}\) , or to \(u^{*}(y) = 0\) . Hence we have \(\mathrm{Q}^{\circ} = \operatorname{Ker}u^{*}\) , and \(\mathrm{Q}\) is the initial subspace of \(u^{*}\) . Since \(u\) is the adjoint of \(u^{*}\) , the final subspace of \(u^{*}\) is also the initial subspace of \(u\) . This proves (i).

The relation \(u(\mathbf{M}) \subset \mathbf{M}\) implies that \(u(\mathbf{M})\) is orthogonal to \(\mathbf{M}^{\circ}\) , and the relation \(u^{*}(\mathbf{M}^{\circ}) \subset \mathbf{M}^{\circ}\) implies that \(u^{*}(\mathbf{M}^{\circ})\) is orthogonal to \(\mathbf{M}\) . But we have \(\langle u(x)|y\rangle = \overline{\langle u^{*}(y)|x\rangle}\) for all \(x \in \mathbf{M}\) and \(y \in \mathbf{M}^{\circ}\) ; hence (ii) follows.

We remark that prop. 4 can be deduced from the general properties of the transpose (II, p.~51, cor. 2) in view of remark 3, V, p.~41.

DEFINITION 3. --- Let E and F be two hilbertian spaces. A mapping \(u \in \mathcal{L}(\mathrm{E};\mathrm{F})\) is said to be partially isometric if \(\|u(x)\| = .\|x\|\) for all x belonging to the initial subspace of u.

Let \(u \in \mathcal{L}(\mathrm{E};\mathrm{F})\) and let \(\mathbf{N}\) be its kernel and \(\mathbf{I}\) its image. To say that \(u\) is partially isometric is the same as saying that the linear mapping \(\tilde{u}: \mathrm{E/N} \to \mathrm{I}\) deduced from \(u\) is isometric (V, p.~13). Then the subspace \(\mathbf{I}\) of \(\mathbf{F}\) is complete, hence closed, and is the final subspace of \(u\) . Consequently, \(u\) induces a hilbertian space isomorphism from the initial subspace of \(u\) onto its final subspace.

PROPOSITION 5. --- Let \(u \in \mathcal{L}(\mathrm{E};\mathrm{F})\) , let \(\mathbf{P}\) be its initial subspace and \(\mathbf{Q}\) be the final subspace. Let \(p\) (resp. \(q\) ) denote the initial (resp. final) orthoprojector of \(u\) . Assume that \(u\) is partially isometric.

\begin{enumerate}
\def\labelenumi{(\roman{enumi})}
\item
  The mapping \(u^{*} \in \mathcal{L}(\mathbf{F};\mathbf{E})\) is partially isometric, with initial subspace \(\mathbf{Q}\) and final subspace \(\mathbf{P}\) . The isomorphism from \(\mathbf{P}\) onto \(\mathbf{Q}\) induced by \(u\) is then the inverse of the isomorphism from \(\mathbf{Q}\) onto \(\mathbf{P}\) induced by \(u^{*}\) .
\item
  We have \(u^{*}u = p\) and \(uu^{*} = q\) .
\end{enumerate}

On account of prop. 4 (i), assertion (i) is a consequence of (ii).

We now prove (ii). Since \(\mathbf{P}\) contains the image of \(u^{*}\) , the mapping \(u^{*}u\) maps \(\mathbf{E}\) into \(\mathbf{P}\) . Let \(x \in \mathbf{E}\) and \(y \in \mathbf{P}\) , then

\[
\langle u ^ {*} u (x) | y \rangle = \langle u (x) | u (y) \rangle .
\]

If x belongs to P, then \(\langle u(x)|u(y)\rangle = \langle x|y\rangle\) by the definition of a partially isometric mapping; if x belongs to the kernel N of u, then \(u(x) = 0\) , hence \(\langle u(x)|u(y)\rangle = 0\) and \(\langle x|y\rangle = 0\) since N and P are orthogonal. Since \(E = P \oplus N\) , we have \(\langle u^{*}u(x) - x|y\rangle = 0\) in all the cases, and so \(u^{*}u\) is the orthoprojector p from E onto P. That \(uu^{*} = q\) follows by interchanging u and \(u^{*}\) in the above.

PROPOSITION 6. --- For every \(u \in \mathcal{L}(\mathrm{E};\mathrm{F})\) , the following conditions are equivalent :

\begin{enumerate}
\def\labelenumi{(\roman{enumi})}
\item
  \(u\) is partially isometric;
\item
  \(u^{*}\) is partially isometric;
\item
  \(u^{*}u\) is an orthoprojector;
\item
  \(uu^{*}\) is an orthoprojector;
\item
  \(uu^{*}u = u;\)
\item
  \(u^{*}uu^{*} = u^{*}\) .
\end{enumerate}

By prop. 5, (i) is equivalent to (ii).

\begin{enumerate}
\def\labelenumi{(\roman{enumi})}
\item
  \(\Rightarrow\) (v): Suppose \(u\) is partially isometric. Then \(u^{*}u\) is the initial orthoprojector of \(u\) by prop. 5. Hence for every \(x \in \mathrm{E}\) , \(u^{*}u(x) - x\) belongs to the kernel of \(u\) , that is, \(uu^{*}u(x) = u(x)\) .
\item
  \(\Rightarrow\) (iii): Suppose that \(uu^{*}u = u\) and let \(p = u^{*}u\) ; then \(p = p^{*}\) and \(p^2 = p\) . Let M (resp. N) be the image (resp. the kernel) of \(p\) . For \(x \in \mathbf{M}\) and \(y \in \mathbf{N}\) , we have \(\langle x|y\rangle = \langle p(x)|y\rangle = \langle x|p^{*}(y)\rangle = \langle x|p(y)\rangle = 0\) . Since M and N are orthogonal, \(p\) is the orthoprojector from E onto M.
\item
  \(\Rightarrow\) (i): Suppose \(p = u^{*}u\) is an orthoprojector with image M and kernel N.
\end{enumerate}

For all \(x\in \mathbf{E}\) , we have

\[
\left\| u (x) \right\| ^ {2} = \langle u ^ {*} u (x) | x \rangle = \langle p (x) | x \rangle .
\]

Hence \(u(x) = 0\) for \(x \in \mathbf{N}\) and \(\| u(x) \| = \| x \|\) for \(x \in \mathbf{M}\) , and so \(u\) is partially isometric with kernel \(\mathbf{N}\) and initial subspace \(\mathbf{M}\) .

We have proved the equivalence of (i), (iii) and (v). Replacing \(u\) by \(u^*\) , we can deduce the equivalence of (ii), (iv) and (vi). This proves prop. 6.
% semantic-fragments: legacy-input-seam
\relax{}
